\section{Disjointness}
\label{sec:disjointness}

\subsection{Blocchi AllDisjointClasses}

L'ontologia dichiara \textbf{12 blocchi \texttt{owl:AllDisjointClasses}}, elencati
nella Tabella~\ref{tab:disjoint}. Ciascun blocco specifica un insieme di classi
mutuamente disgiunte: nessuna istanza pu\`o appartenere a pi\`u di una classe
nello stesso blocco.

\begin{table}[h]
\centering
\caption{Blocchi AllDisjointClasses}
\label{tab:disjoint}
\begin{tabular}{ll}
\toprule
\textbf{\#} & \textbf{Classi} \\
\midrule
1 & VideoGame, Developer, Publisher, Genre, Platform, Character, Franchise, Award, \\
  & GameEngine, GameMode, Review, GameEvent \\
2 & ConsolePlatform, PCPlatform, MobilePlatform \\
3 & PlayStationPlatform, XboxPlatform, NintendoPlatform \\
4 & WindowsPlatform, MacOSPlatform, LinuxPlatform \\
5 & IOSPlatform, AndroidPlatform \\
6 & IndustryAward, EditorialAward, UserAward \\
7 & PlayerCharacter, NonPlayerCharacter \\
8 & ProprietaryEngine, OpenSourceEngine \\
9 & ReleasedGame, EarlyAccessGame, CancelledGame, DelistedGame \\
10& IndieGame, AAAGame \\
11& SinglePlayerMode, MultiplayerMode, CoopMode, MMOMode \\
12& AnnouncementEvent, ReleaseEvent, UpdateEvent, DelistingEvent \\
\bottomrule
\end{tabular}
\end{table}

\subsection{Importanza delle Disjointness}

Le dichiarazioni di disjointness svolgono due ruoli fondamentali:

\begin{enumerate}[nosep]
  \item \textbf{Rilevazione di inconsistenze nei dati Wikidata}. Se un'entit\`a
    Wikidata venisse erroneamente classificata come, ad esempio, sia
    \texttt{ConsolePlatform} che \texttt{PCPlatform}, un reasoner DL completo
    (HermiT, Pellet) rileverebbe l'inconsistenza. Questo meccanismo funge da
    \textit{sanity check} automatico sulla qualit\`a dei dati;
  \item \textbf{Ragionamento DL}. Le disjointness abilitano inferenze per
    contraddizione. Ad esempio, sapendo che \texttt{PlayerCharacter} e
    \texttt{NonPlayerCharacter} sono disgiunte, se un personaggio viene
    classificato come \texttt{PlayerCharacter}, si pu\`o inferire che
    \textbf{non} \`e \texttt{NonPlayerCharacter}.
\end{enumerate}

\subsection{Nota su SinglePlayerGame e MultiplayerGame}

Si noti che \texttt{SinglePlayerGame} e \texttt{MultiplayerGame} \textbf{non}
compaiono in alcun blocco di disjointness. La scelta \`e deliberata: nel dominio
dei videogiochi, \`e comune che un gioco supporti sia la modalit\`a single-player
che multiplayer (es. Grand Theft Auto V, Elden Ring, Call of Duty). Dichiararle
disgiunte produrrebbe inconsistenze su una porzione significativa delle istanze.

Al contrario, le modalit\`a \texttt{SinglePlayerMode}, \texttt{MultiplayerMode},
\texttt{CoopMode} e \texttt{MMOMode} (blocco \#11) \textbf{sono} disgiunte:
una specifica modalit\`a di gioco \`e esattamente una di queste quattro.
